\documentclass[10pt]{article}
\usepackage[margin=0.72in]{geometry}
\usepackage{booktabs}
\usepackage{array}
\usepackage{tabularx}
\usepackage[T1]{fontenc}

\setlength{\emergencystretch}{2em}
\newcolumntype{L}[1]{>{\raggedright\arraybackslash}p{#1}}
\newcolumntype{Y}{>{\raggedright\arraybackslash}X}

\title{SM120 Sparse MMA Instruction Diagnostic}
\author{}
\date{Measurement date: 2026-05-09}

\begin{document}
\maketitle

This directory is now a diagnostic appendix. Do not use this report to state
the sparse-vs-dense speedup for a matrix workload.

The primary performance report is in the sibling
\texttt{sm120\_mma\_workload\_sweep} directory. Use its Markdown and PDF reports
for speedup conclusions. That report compares sparse and dense at the same
\texttt{A[M,K] x B[K,N]} workload.

\section{Why This Is Diagnostic Only}

The original benchmark in this directory times individual warp-level MMA packet
forms. It is useful for checking whether RTX 5090 / \texttt{sm\_120} accepts and
runs a PTX instruction spelling, but its aggregate medians are not a clean
workload performance metric.

The confusing case was F16/F16 ordered sparse:

\begin{itemize}
\item \texttt{m16n8k16} sparse compared with one dense \texttt{m16n8k16} packet
is near parity.
\item \texttt{m16n8k32} sparse compared with two dense \texttt{m16n8k16}
packets is near 2x.
\item Taking a median across those two packet classes gives about 1.5x, but
that is not the speedup of a same-size matrix workload.
\end{itemize}

The same-workload benchmark fixes this by choosing one matrix shape, giving
sparse and dense the same \texttt{M}, \texttt{N}, and \texttt{K}, and timing the
tensor-core work for that equal workload.

\section{Current Same-Workload Result To Quote}

From \texttt{sm120\_mma\_workload\_sweep/REPORT.md}:

\begin{center}
\small
\begin{tabular}{lrr}
\toprule
Mode & Overall speedup & Overall cycle ratio \\
\midrule
Throughput & 1.946543x & 0.513732x \\
Latency & 1.962671x & 0.509509x \\
\bottomrule
\end{tabular}
\end{center}

For F16 inputs with F16 accumulation:

\begin{center}
\small
\begin{tabular}{lrr}
\toprule
Mode & F16/F16 speedup & F16/F16 cycle ratio \\
\midrule
Throughput & 2.039617x & 0.490292x \\
Latency & 2.037302x & 0.490849x \\
\bottomrule
\end{tabular}
\end{center}

\section{What This Diagnostic Still Covers}

This diagnostic remains useful for:

\begin{itemize}
\item verifying that baseline \texttt{sm\_120} accepts warp-level
\texttt{mma.sp} and \texttt{mma.sp::ordered\_metadata} forms;
\item identifying that \texttt{tcgen05.*} is not accepted for
\texttt{.target sm\_120};
\item checking compiler advisories for plain \texttt{.sp};
\item checking whether a sparse PTX spelling runs without launch/runtime
failure;
\item preserving exact sparse and dense PTX spellings in
\texttt{results/sm120\_mma\_instruction\_pairs.csv}.
\end{itemize}

\section{Coverage}

The diagnostic sweep completed 176 sparse/dense packet rows with no build or
runtime failures on an NVIDIA GeForce RTX 5090.

\begin{center}
\small
\begin{tabularx}{\linewidth}{L{0.16\linewidth}L{0.25\linewidth}Y L{0.14\linewidth}}
\toprule
Family & Sparse packet shapes & Input formats & Accum. \\
\midrule
F16/F32 & \texttt{m16n8k16}, \texttt{m16n8k32} & \texttt{f16 x f16} & \texttt{f32} \\
F16/F16 & \texttt{m16n8k16}, \texttt{m16n8k32} & \texttt{f16 x f16} & \texttt{f16} \\
BF16 & \texttt{m16n8k16}, \texttt{m16n8k32} & \texttt{bf16 x bf16} & \texttt{f32} \\
TF32 & \texttt{m16n8k8}, \texttt{m16n8k16} & \texttt{tf32 x tf32} & \texttt{f32} \\
FP8 & \texttt{m16n8k64} & \texttt{e4m3/e5m2 x e4m3/e5m2} & \texttt{f32} \\
INT8 & \texttt{m16n8k32}, \texttt{m16n8k64} & \texttt{u8/s8 x u8/s8} & \texttt{s32} \\
INT4 & \texttt{m16n8k64}, \texttt{m16n8k128} & \texttt{u4/s4 x u4/s4} & \texttt{s32} \\
\bottomrule
\end{tabularx}
\end{center}

For INT8 and INT4, both normal and \texttt{.satfinite} accumulation are
covered. For sparse instructions, both plain \texttt{.sp} and ordered metadata
are covered.

\section{Diagnostic Findings}

\begin{itemize}
\item The useful sparse spelling on RTX 5090 is
\texttt{mma.sp::ordered\_metadata.sync.aligned}.
\item Plain \texttt{mma.sp.sync.aligned} is not a good performance path for this
machine. \texttt{ptxas} emits a performance advisory, and the measured packet
path is poor for F16/F16 and INT8.
\item INT4 sparse packet forms are slower than their dense peers in both this
diagnostic and the same-workload sweep.
\item \texttt{tcgen05.*} is not supported by \texttt{ptxas} for
\texttt{.target sm\_120}; those forms remain outside the RTX 5090 baseline path.
\end{itemize}

\section{Artifacts}

\begin{center}
\small
\begin{tabularx}{\linewidth}{L{0.40\linewidth}Y}
\toprule
Path & Description \\
\midrule
\texttt{sm120\_mma\_sparse.cu} & Instruction-packet CUDA diagnostic \\
\texttt{benchmark.py} & Builds/runs sparse and dense packet pairs \\
\texttt{summarize\_results.py} & Generates diagnostic CSV summaries \\
\texttt{results/sm120\_mma\_full.csv} & Raw 176-row diagnostic sweep \\
\texttt{results/sm120\_mma\_instruction\_pairs.csv} & Exact sparse/dense PTX spelling pairs \\
\texttt{results/sm120\_mma\_overall.csv} & Packet-level aggregate summary \\
\texttt{results/sm120\_mma\_summary\_by\_family.csv} & Packet-level family summary \\
\texttt{results/sm120\_mma\_extremes.csv} & Packet-level best/worst rows \\
\bottomrule
\end{tabularx}
\end{center}

\section{Reproduction}

\begin{verbatim}
cd sm120_mma_sparse
python3 benchmark.py --mode both --preset full \
    --sparse-variant both --output results/sm120_mma_full.csv --overwrite
python3 summarize_results.py results/sm120_mma_full.csv --out-dir results
make report
\end{verbatim}

Use the same-workload report for performance conclusions:

\begin{verbatim}
cd ../sm120_mma_workload_sweep
python3 workload_benchmark.py \
    --mode both --format-preset full --size-preset full \
    --sparse-variant ordered --target-packets-per-warp 4096 \
    --warmup-launches 3 \
    --output results/sm120_mma_workload_sweep.csv --overwrite
python3 summarize_workload.py results/sm120_mma_workload_sweep.csv --out-dir results
make report
\end{verbatim}

\end{document}
